理想 Chaplygin 雪橇在水平平面运动：刚体质量 \(m>0\)、质心 \(C\)、质心转动惯量 \(I>0\)。刀刃接触点为 \(A=C-a\mathbf e_1\)，\(a>0\)，其中 \(\mathbf e_1=(\cos\theta,\sin\theta)\)、\(\mathbf e_2=\hat{\mathbf z}\times\mathbf e_1\)。刀刃禁止横向滑动，故 \(\mathbf v_A=u\mathbf e_1\)，但允许沿轴自由运动；刀刃反力仅沿 \(\mathbf e_2\)，无其他水平外力。令 \(\omega=\dot\theta\)、\(J=I+ma^2\)，总动能 \(E>0\)。
(1) 从质心速度与牛顿平动、角动量方程推导非完整动力学 \(\dot u,\dot\omega\) 及刀刃约束反力，求出动能的二次型，并说明不能把刀刃条件误当可积位置约束。
(2) 证明能量守恒；对非直线的 \(\omega\ne0\) 轨道，以 \(v=\sqrt{2E/m}\) 为参数求双曲函数精确解，逐项代回动力学方程并识别特殊的直线运动解。
(3) 求精确朝向 \(\theta(t)\) 及整条异宿轨道的总转角，线性化判定 \(u=\pm v,\omega=0\) 的稳定性，并解释 \(a\to0\) 时为什么不能直接交换长时间极限。
(4) 由 \(\dot{\mathbf r}_A=u\mathbf e_1\) 在 \(u(t_0)=0\) 的时刻展开轨迹至三阶，证明半三次尖点结构；同时计算该时刻的约束反力、约束功率与量纲。
